\documentclass{article}
\usepackage[T1]{fontenc}
\usepackage[utf8]{inputenc}
\usepackage{lmodern}
\usepackage[polish,shorthands=off]{babel}
\usepackage{amsmath,mathtools,amsfonts,stackengine}
\usepackage{pgf-umlsd}
\stackMath
\usepackage{amssymb}
\title{Zadanie 6 z listy 1}
\author{Jakub Kowal}
\date{}

\begin{document}
\maketitle
\section*{Schemat}

\begin{center}
    

\begin{sequencediagram}
\newinst{D}{Dowodzący}
\newinst{W}{Weryfikator}

\mess{D}{$X = g^{x}$}{W}
\mess{W}{$c$}{D}
\mess{D}{$s = \lparen a + cx\rparen \lparen c + ax\rparen$}{W}
\mess{D}{$W=g^{ax^{2}}$}{W}
\mess{D}{$Z=g^{xa^{2}}$}{W}
\begin{messcall}{W}{$\text{ACC jeśli } g^{s}=(A X^{c}W)^{c}Z$}{D}
\end{messcall}
\end{sequencediagram}
\end{center}

\section{}
\begin{center}
    $g^{s} = {(AX^{c}W)}^{c}Z$\\
    $g^{(a + cx)(c + ax)}=(g^{a}g^{cx}g^{ax^{2}})^{c}g^{xa^{2}}$\\
    $g^{ac + a^{2}x + c^{2}x + acx^{2}} = g^{ac}g^{c^{2}x}g^{acx^{2}}g^{xa^{2}}$\\
    $g^{ac + a^{2}x + c^{2}x + acx^{2}} = g^{ac + a^{2}x + c^{2}x + acx^{2}}$
\end{center}

\section{}
Nie wymusza znajmości sekretu a, co będzie udowodnione w nastepnym punkcie.
\section{}
Weryfikator akceptuje jeśli spełnione jest $g^{s}=(A X^{c}W)^{c}Z$. Możemy wyprowadzić:
\begin{center}
    $Z=\frac{g^{s}}{{(AX^{c}W)}^{c}}$
\end{center}Dowodzący dosyła s, W, Z, więc wysyłając $s = 0$ i $W = 1$ otrzymujemy równanie:
\begin{center}
    $Z=\frac{1}{{(AX^{c})}^{c}}$
\end{center}
Znając x oraz A obliczamy Z, które spełni warunki akceptacji.

\section{}
W tym przypadku wysoka złożoność wzrorów jedynie zaciemnia obraz, bo jak udowodniłem łatwo da się obejść zabezpieczenie.

\section{}
Skuteczność kryptografii nie polega na tym, żeby równania były trudne do odczytania, tylko na tym, żeby atakujący nie miał żadnej swobody ruchu. Bezpieczne protokoły (takie jak klasyczny Schnorr) są proste i eleganckie – bezpieczeństwo bierze się z tego, że napastnik musi rozwiązać konkretny trudny problem. Dodawanie dużej ilości dodatkowych parametrów, może powodować więcej miejsc na luki w systemie.

\end{document}
